% title:          Decimal digits of powers of 5+sqrt(26)
% area:           digits-bases, irrationality
% methods:        algebraic-conjugates
% difficulty:
% difficulty-ai:  easy
% status:
% review:         none
% --- history: all optional, leave EMPTY if unknown ---
% origin:         folklore
% origin-date:
% origin-number:
% also-in:        course
% taken-from:
% references:     Ravi Vakil, Stanford Putnam seminar, Problem Solving Masterclass Week 4 (Monday 17 Nov 2003), Problem 4, presented by 'Andy', no source given; same wording as the bank statement - http://math.stanford.edu/~vakil/putnam03/03mc4.pdf ; seminar page listing the masterclass handouts - http://math.stanford.edu/~vakil/putnam03/
% history:        ai
\begin{problem}
Prove that the decimal part of
\[
{(5 + \sqrt{26})}^{n}
\]
begins with either $n$ zeros or $n$ nines for all positive integers $n$.
\end{problem}
\begin{solution}
It is clear that ${(5 + \sqrt{26})}^{n} + {(5 - \sqrt{26})}^{n}$ is an integer, whether by expanding the products and noting that all odd powers of $\sqrt{26}$ cancel, or by using Galois theory - sum of algebraic integers is algebraic integer, sum of all roots is a rational, rational algebraic integer is integer.\\
What remains is to show $|{(5 - \sqrt{26})}^{n}| < 10^{-n}$, or simply $5 - \sqrt{26} = \frac{1}{5 + \sqrt{26}} < \frac{1}{5+5} < \frac{1}{10}$.\\

Note: These are called Pisot-Vijayaraghavan numbers.
\end{solution}
